

\subsection{Interface}

RetailBench models single-store supermarket operations as a partially
observable decision process over 180 simulated days. Agents manage
pricing, replenishment, and supplier selection through tool-based
interactions, as illustrated in Figure~\ref{fig:retailbench_framework}.
These decisions are interdependent and have delayed consequences:
pricing and replenishment affect future sales and inventory, while
delivery delays, product expiration, and market changes influence
subsequent operations and financial outcomes.

Formally, we model RetailBench as a partially observable Markov
decision process (POMDP):
\begin{equation}
\label{eq:pomdp}
\mathcal{M} = (\mathcal{S}, \mathcal{A}, \mathcal{O}, T, \Omega, R),
\end{equation}
where $\mathcal{S}$ is the state space, $\mathcal{A}$ the tool-action
space, $\mathcal{O}$ the observation space, $T$ the transition function,
$\Omega$ the observation function, and $R$ the business-outcome
function. Agents observe the environment only through available tools,
without access to its full underlying state. Appendix~\ref{app:state_action_spaces}
provides the detailed state and action spaces.

At day $t$ and intra-day step $k$, the environment state is
factorized as:
\begin{equation}
\label{eq:latent_state}
S_{t,k} = (P, I, V, C, D, E, F)_{t,k},
\end{equation}
where $P$ denotes product attributes and prices, $I$ inventory and
pending orders, $V$ shelf assortment, $C$ supplier conditions,
$D$ demand factors, $E$ external events, and $F$ financial state.
These components jointly characterize retail operations under
inventory, shelf-capacity, and financial constraints.

Agents interact with RetailBench through three categories of tools:
read-only tools for information acquisition,
auxiliary tools for memory and computation, and
operational tools for business decisions.
At day $t$ and interaction step $k$, let $S_{t,k}$ denote the
environment state, $h_{t,k}$ the agent's interaction history,
and $a_{t,k}$ the selected tool action.
The interaction is formalized as
\begin{equation}
\label{eq:agent_interaction}
\begin{aligned}
\mathcal{A} &=
\mathcal{A}^{\mathrm{read}}
\cup \mathcal{A}^{\mathrm{other}}
\cup \mathcal{A}^{\mathrm{act}},\\
o_{t,k} &= \Omega(S_{t,k},a_{t,k}),
\, a_{t,k}\in\mathcal{A}^{\mathrm{read}},\\
z_{t,k} &= G(h_{t,k},a_{t,k}),
\, a_{t,k}\in\mathcal{A}^{\mathrm{other}},\\
S_{t,k+1} &= T(S_{t,k},a_{t,k}),
\, a_{t,k}\in
\mathcal{A}^{\mathrm{act}}\!\setminus\!\{a^{\mathrm{end}}\}.
\end{aligned}
\end{equation}

Here, $\Omega$ is the observation function that returns
environment information $o_{t,k}$, while $G$ represents
auxiliary operations producing outputs $z_{t,k}$ from
the agent's interaction history.
The transition function $T$ updates the environment state
following operational actions, whereas read-only and
auxiliary tools leave the simulator state unchanged.
Specifically, read-only tools expose inventory, sales,
supplier quotes, and customer feedback without revealing
latent demand or supplier-quality parameters.
Auxiliary tools support memory and computation, while
operational tools enable pricing, replenishment,
shelf management, and persistent note updates.

Time advances only when the agent invokes \texttt{end\_today},
which triggers the daily settlement:
\begin{equation}
\label{eq:end_today_transition}
S_{t+1,0}
= T(S_{t,k},a^{\mathrm{end}})
\equiv T_{\mathrm{day}}(S_{t,k}).
\end{equation}
This transition processes deliveries, sales, customer feedback,
returns, product expiration, operating costs, and market changes.
Consequently, current decisions shape future inventory, cash flow,
and operational opportunities, requiring agents to balance
short-term performance with long-term profitability and survival.

\subsection{Data-Grounded Simulator}

RetailBench integrates real-world retail data with controlled simulation
mechanisms to model long-horizon business operations. Its data sources
provide empirical foundations for retail demand, pricing, customer
feedback, and external events.

We use the Dominick's dataset~\citep{DominicksDataset}, which contains
store-level records of product categories, sales, prices, and
cost-related information. These records inform the product catalog
and empirical priors for demand, store traffic, pricing, and
procurement costs. We select 96 products across 20 categories
with sufficient records to estimate product-level demand and
category-level substitution patterns.

Customer feedback is grounded in Amazon Reviews
2023~\citep{hou2024bridging}, which provides category-specific review
pools organized by ratings, sentiment, and product aspects.
Realized sales generate reviews whose ratings depend on supplier
quality, with comments sampled from the corresponding generated reviews pool.
External events are constructed from financial news
articles~\citep{ashraq2025financialNewsArticles}, covering neutral,
category-level, and product-level events.
Agents observe news and historical feedback, while latent event
parameters govern their effects on demand and supply.

Together, these sources provide data-grounded operational signals
within a controlled, partially observable simulation environment.


\subsection{Modules}
\label{sec:module_logic}

RetailBench comprises seven interconnected modules: Inventory,
Suppliers \& Orders, Shelf, Customer Demand, Reviews \& Returns,
News Events, and Finance \& Records
(Appendix~\ref{app:simulator_modules}).
Together, these modules govern inventory availability, supplier
conditions, demand realization, customer feedback, external events,
and financial outcomes.

Module interactions introduce dependencies across operational decisions
and time. Replenishment affects future inventory through delivery
delays, while pricing and shelf allocation influence realized sales.
Inventory aging, returns, and operating costs further affect
subsequent inventory availability and financial conditions.
Consequently, operational decisions have delayed effects that
shape future states and constrain subsequent actions.

Detailed module definitions and day-level update equations are
provided in Appendix~\ref{app:simulator_modules}
and Appendix~\ref{app:simulator_update_equations}.


\subsection{Scope}

RetailBench focuses on integrated, single-store retail operations
rather than isolated operational subproblems
(Section~\ref{sec:related_work}; Table~\ref{tab:benchmark_axes}).
By integrating data-grounded business conditions, interdependent decisions, and delayed consequences, RetailBench provides a controlled yet realistic retail environment for evaluating long-horizon agent decision-making under partial observability.